% Current-cohort replacements for the user's supplied paragraphs.
% Requires booktabs, multirow, graphicx, xcolor and the manuscript's \fix command.
% REVIEW REQUIRED: unchanged prose conflicts with some current measurements.
% See REVIEW.txt before using this fragment in a manuscript.
% Intro cohort: measured OuterYan pairs plus verified native-plan fallbacks.
% Excludes selected-but-untimed rows, timeouts and pending queries.
Over @@intro_n@@ queries, DuckOuterYan picks a rewritten plan for @@selected_n@@ of them and is faster than DuckDB on these by @@intro_gm@@ in geometric mean and by up to @@intro_max@@, while the remaining queries run with DuckDB's native plan (Section~\ref{sec:experiments}).

% Reporting corpus after the explicit speed-up screen, before timing-validity filtering.
@@benchmark_sentence@@

@@cohort_note@@

@@runtime_table@@

\subsubsection{Overall Performance Analysis}

Across the @@selected_n@@ valid measured query instances that apply the OuterYan plan, OuterYan achieves a @@overall_gm@@ geometric-mean speed-up over DuckDB. The geometric-mean speed-ups for Graph, TPC-H, TPC-DS, and JOB are @@graph_gm@@, @@tpch_gm@@, @@tpcds_gm@@, @@job_gm@@, respectively. The maximum per-query speed-ups for these workloads are @@graph_max@@, @@tpch_max@@, @@tpcds_max@@, @@job_max@@, respectively, with a maximum of @@overall_max@@ across these OuterYan-selected instances.

Their total-time ratio is @@total_ratio@@. We exclude timed-out queries and report performance statistics only for valid measured queries for which an OuterYan plan is selected. Overall, the selected OuterYan plans achieve substantial speedups over DuckDB across all four benchmarks.

\subsubsection{Case Query Analysis}

\fix{TODO}

\subsubsection{Execution Breakdown}
\label{sec:phase-breakdown}

\begin{figure*}[t]
\centering\includegraphics[width=\textwidth]
{figures/phase_breakdown_outeryan.pdf}
\caption{Execution breakdown for representative OuterYan-selected queries.}
\label{fig:phase-breakdown}
\end{figure*}

Figure~\ref{fig:phase-breakdown} shows representative queries executed with an OuterYan plan. The left panel breaks the profiled CPU work into Round~1 bottom-up semi-joins, Round~2 top-down semi-joins, Round~3 real joins, and the remaining operators. Each bar is normalized to 100\%, and the value at the end reports the total profiled CPU work. The right panel compares DuckDB and OuterYan running times for the same queries on a logarithmic scale. The relative marker positions indicate whether OuterYan is faster or slower, and the value on the right gives the corresponding speed-up.

\begin{table}[t]
\centering\small\setlength{\tabcolsep}{5pt}
\caption{Phase-time breakdown of OuterYan plans.}
\label{tab:phase-by-benchmark}
\begin{tabular}{lrrrrr}
\toprule
Family & $n$ & R1 CPU & R2 CPU & R3 CPU & Other CPU \\
\midrule
@@phase_rows@@
\bottomrule
\end{tabular}
\end{table}

The execution breakdown differs across workloads. For the Graph dataset, most execution work remains in Round~3 because graph-style many-to-many joins can produce large intermediate results. In contrast, joins in TPC-H, TPC-DS, and JOB are often constrained by PK--FK, which limit intermediate result sizes and make the cost of the two semi-join rounds more substantial relative to subsequent joins. Overall, OuterYan is beneficial when the reduction in Round~3 join work is sufficient to compensate for the additional semi-join overhead. In cases such as H-RD-03, the semi-join reductions provide limited pruning while adding extra work, resulting in a modest slowdown.

\subsubsection{Optimization Time}
\label{sec:optimization-time}
Because the DuckOuterYan pipeline is implemented within DuckDB's optimizer, we separately measure the total time spent on all OuterYan optimization steps and report its fraction of the overall query running time.
Table~\ref{tab:optimization-time} shows that OuterYan's optimization ratio remains low, staying below \fix{1\%} for @@optimization_below@@ and reaching at most @@optimization_max@@. Although OuterYan's optimization time is often higher than that of native DuckDB, the additional optimization overhead remains negligible compared with the execution-time savings achieved by its plans.

\begin{table*}[t]
\centering\scriptsize\setlength{\tabcolsep}{3.2pt}
\caption{Running time, query statistics, and optimization overhead for representative queries. Optimization ratio is $100\,T_{\rm opt}/T_{\rm runtime}$ for each configuration, using separate optimizer profiles and unprofiled runtime medians. Attributes count distinct referenced columns per relation alias, including expanded SELECT stars, join predicates and filters. Queries are selected by family, shape and DuckDB runtime ranks, without ranking speed-up.}
\label{tab:optimization-time}
\resizebox{\textwidth}{!}{%
\begin{tabular}{@@optimization_columns@@}
\toprule
@@optimization_header@@
\midrule
@@optimization_rows@@
\bottomrule
\end{tabular}}
\end{table*}

\subsubsection{PostgreSQL Plan Comparison}
\label{sec:pgplan-comparison}

Building on the motivating example in Section~\ref{sec:intro}, we examine two existing benchmark variants: a five-table, right-deep JOB-based query (SQLStorm 13193)~\cite{job} and a six-table, bushy TPC-H Q8 variant at SF100~\cite{tpch}. Both mix inner and outer joins. We retain their benchmark predicates and execute them without ORDER BY or LIMIT. PG-Plan denotes PostgreSQL~16.2's selected logical join tree, including join types, replayed in DuckDB; DuckDB and OuterYan denote execution with OuterYan disabled and enabled, respectively. All modes use the same DuckDB Release build, one thread, and a 4\,GB memory limit. We verify complete result equivalence and report median runtimes over five runs in Table~\ref{tab:pgplan-runtime}.

For the JOB query, both baselines join \texttt{cast\_info} with \texttt{aka\_name} early. PostgreSQL's enumeration trace rejects joining \texttt{aka\_title} with \texttt{cast\_info} first because of an outer-join constraint. OuterYan uses semijoin reduction to preserve the required matches while enabling this grouping and deferring the join with \texttt{aka\_name}. This reduces the largest join output from \fix{34.36 million} rows under PG-Plan and \fix{45.37 million} under DuckDB to \fix{0.47 million}. OuterYan consequently achieves speedups of @@pg_job_speedup@@ and @@native_job_speedup@@, respectively.

For the TPC-H query, PG-Plan joins \texttt{orders} with \texttt{lineitem} before \texttt{customer}, whereas OuterYan first combines \texttt{orders} with the customer branch. Both apply the order-date filter early, but OuterYan additionally uses semijoin reduction to shrink subsequent join inputs: only \fix{18.45 million} of the \fix{600.04 million} lineitem rows reach the main joins. Native DuckDB applies the date filter after large intermediate joins and incurs temporary disk I/O. OuterYan is therefore @@pg_tpch_speedup@@ faster than PG-Plan and @@native_tpch_speedup@@ faster than DuckDB. The JOB case demonstrates a useful semijoin-enabled grouping beyond PostgreSQL's enumeration rules, while both cases show the performance benefits of combining reordering with reduction.

\begin{table}[t]
  \centering
  \small
  \setlength{\tabcolsep}{5pt}
  \caption{Median query running time in DuckDB (ms). PG-Plan replays PostgreSQL's selected join tree.}
  \label{tab:pgplan-runtime}
  \begin{tabular}{lrrr}
    \toprule
    Query & PG-Plan & DuckDB & OuterYan \\
    \midrule
    JOB 13193\textsuperscript{a}
      & @@pg_job_ms@@ & @@native_job_ms@@ & @@oy_job_ms@@ \\
    TPC-H Q8 / SF100\textsuperscript{b}
      & @@pg_tpch_ms@@ & @@native_tpch_ms@@ & @@oy_tpch_ms@@ \\
    \bottomrule
  \end{tabular}
  \par\smallskip
  \begin{minipage}{\linewidth}\footnotesize
  @@pg_footnotes@@
  \end{minipage}
\end{table}
